%%%%%%%%%%%%%%%%%%%%%%%%%%%%%%%%%%%%%%%%%%%%%%%%%%%%%%%%%%%%%%%%%%%%%%%%%%%%%%%
% SECTION 11.10: Temporal Dynamics
% Axioms: AWX-A56 to AWX-A65
% Appendix AWA Reference: Section 3.8
%%%%%%%%%%%%%%%%%%%%%%%%%%%%%%%%%%%%%%%%%%%%%%%%%%%%%%%%%%%%%%%%%%%%%%%%%%%%%%%

Awareness is not static but evolves over time through decay, fluctuation, 
habituation, and life events. This section specifies the temporal dynamics 
that govern how awareness changes before, during, and after the Moment of Truth.

\subsubsection{Situational Variation (AWX-A56 to A57)}

\paragraph{AWX-A56: Situativity.}
Awareness varies dramatically across situations, even for the same person 
considering the same consequences:
\begin{equation}
A(t, s) \neq A(t, s') \quad \text{for situations } s \neq s'
\end{equation}

\paragraph{AWX-A57: Within-Situation Fluctuation.}
Even within a single situation, awareness fluctuates moment-to-moment:
\begin{equation}
A(t) = \bar{A} + \epsilon(t), \quad \epsilon \sim N(0, \sigma^2_A)
\end{equation}
This stochastic component explains why identical situations can produce 
different behaviors.

\subsubsection{The Decay Function (AWX-A58)}

\begin{axiom}[AWX-A58: Temporal Decay]
Awareness of future consequences decays exponentially with temporal distance:
\begin{equation}
A(t) = A_0 \cdot e^{-\lambda t}
\label{eq:temporal-decay}
\end{equation}
where $\lambda$ is the decay rate and $t$ is the time until consequences occur.
\end{axiom}

The decay rate $\lambda$ varies by:
\begin{itemize}
    \item \textbf{Utility type}: $\lambda_{INU} > \lambda_{KNU} > \lambda_{IDN}$
    \item \textbf{Individual differences}: Higher for present-biased individuals
    \item \textbf{Context}: Higher under stress, lower with institutional support
\end{itemize}

\begin{example}[Retirement Savings (CASE-03)]
Consider saving for retirement 30 years away:
\begin{itemize}
    \item Initial awareness: $A_0 = 0.05$ (low baseline)
    \item Decay rate: $\lambda = 0.12$ (typical for financial outcomes)
    \item Time horizon: $t = 30$ years
\end{itemize}
\[
A_{retirement} = 0.05 \cdot e^{-0.12 \times 30} = 0.05 \cdot e^{-3.6} = 0.05 \cdot 0.027 = 0.00137
\]
A \textbf{657:1 ratio} between present consumption awareness ($A \approx 0.9$) 
and future retirement awareness ($A \approx 0.00137$).
\end{example}

\subsubsection{Reactivation and Habituation (AWX-A59 to A60)}

\paragraph{AWX-A59: Reactivation.}
Dormant awareness can be reactivated by situational triggers:
\begin{equation}
A_{reactivated} = A_{dormant} \cdot (1 + \rho \cdot Trigger)
\end{equation}
where $\rho$ measures trigger effectiveness.

\paragraph{AWX-A60: Habitual Attenuation.}
Repeated exposure to the same information reduces its awareness-generating 
capacity:
\begin{equation}
A_{n} = A_1 \cdot \eta^{n-1}, \quad \eta < 1
\end{equation}
where $n$ is the number of exposures and $\eta$ is the habituation factor.

This explains why repeated health warnings lose effectiveness over time: 
the 100th cigarette warning generates less awareness than the first.

\subsubsection{Life Course Dynamics (AWX-A61 to A62)}

\paragraph{AWX-A61: Life Phase Modulation.}
Awareness patterns vary systematically across life phases:

\begin{table}[htbp]
\centering
\caption{Awareness Patterns Across Life Phases}
\label{tab:life-phases}
\begin{tabular}{@{}lll@{}}
\toprule
\textbf{Phase} & \textbf{Dominant $A$} & \textbf{Reduced $A$} \\
\midrule
Adolescence & $A_{IDN}$, $A^{comm}$ & $A_{long-term}$ \\
Early Career & $A_{INU}$, $A_{financial}$ & $A_{health}$ \\
Parenthood & $A^{comm}$, $A_{KNU}$ & $A^{self}_{INU}$ \\
Late Career & $A_{legacy}$, $A_{IDN}$ & $A_{novelty}$ \\
Retirement & $A_{health}$, $A_{social}$ & $A_{financial}$ \\
\bottomrule
\end{tabular}
\end{table}

\paragraph{AWX-A62: Life Events.}
Major life events can permanently shift awareness baselines:
\begin{equation}
A_{post-event} = A_{pre-event} + \Delta A_{event}
\end{equation}

Examples of awareness-shifting events:
\begin{itemize}
    \item Parenthood: Increases $A_{KNU}$ and $A_{long-term}$
    \item Health crisis: Increases $A_{health}$
    \item Job loss: Increases $A_{financial}$ and $A_{uncertainty}$
    \item Bereavement: Increases $A_{mortality}$ and $A_{relationships}$
\end{itemize}

\subsubsection{Shock and Stabilization (AWX-A63 to A64)}

\paragraph{AWX-A63: Shock Triggers.}
Sudden, unexpected events create temporary awareness spikes:
\begin{equation}
A_{shock}(t) = A_{base} + \Delta A_{shock} \cdot e^{-\mu(t-t_{shock})}
\end{equation}
where $\mu$ governs how quickly the shock effect dissipates.

\paragraph{AWX-A64: Stabilization.}
After disruption, awareness tends toward a new stable equilibrium:
\begin{equation}
\frac{dA}{dt} = -\alpha(A - A_{equilibrium})
\end{equation}

This explains post-crisis behavior: initial panic (high $A$) followed by 
gradual return to baseline, but potentially at a new equilibrium level.

\subsubsection{Journey Integration (AWX-A65)}

\paragraph{AWX-A65: Touchpoint Accumulation.}
In customer or citizen journeys, awareness accumulates across touchpoints:
\begin{equation}
A_{journey} = 1 - \prod_{i=1}^{n}(1 - A_i \cdot w_i)
\end{equation}
where $A_i$ is awareness at touchpoint $i$ and $w_i$ is the touchpoint weight.

This has implications for behavioral intervention design:
\begin{itemize}
    \item Multiple touchpoints are more effective than single interventions
    \item Early touchpoints set the frame for later ones
    \item Consistency across touchpoints reinforces awareness
    \item Contradictory touchpoints create confusion and reduce net awareness
\end{itemize}

\subsubsection{Temporal Implications for Intervention Design}

The temporal dynamics suggest several intervention principles:

\begin{enumerate}
    \item \textbf{Timing matters}: Intervene close to the Moment of Truth 
          to minimize decay
    \item \textbf{Refresh regularly}: Combat habituation with varied messaging
    \item \textbf{Leverage life events}: Major transitions are windows for 
          awareness change
    \item \textbf{Design journeys}: Multiple coordinated touchpoints 
          outperform single interventions
    \item \textbf{Anticipate shock decay}: Post-crisis interventions must 
          account for rapid return to baseline
\end{enumerate}
